\section*{Capítulo \ref{ch:b3:quantitative-geomorphology} --- Geomorfologia quantitativa}

\begin{solution}{exo:b3:quantitative-geomorphology:1}
$a/\sqrt{\pi\kappa t} = 0.30 - 0.04 = 0.26$, logo $\kappa t = (2/0.26)^2/\pi =
\qty{18.8}{m^2}$, e $t = 18.8/(5\times10^{-3}) = \qty{3.8}{kyr}$, dentro de um
fator de dois para mais ou para menos, dado o intervalo de $\kappa$.
\end{solution}

\begin{solution}{exo:b3:quantitative-geomorphology:2}
$E = \kappa\,|\nabla^2 z| = 4\times10^{-3}\times0.012 =
\qty{4.8e-5}{m/yr}$, \qty{0.048}{mm/yr}. Relevo $E L^2/(2\kappa) =
0.012\times80^2/2 = \qty{38.4}{m}$.
\end{solution}

\begin{solution}{exo:b3:quantitative-geomorphology:3}
$E = 2\times10^{-6}\times(25\times10^6)^{0.5}\times0.01 = 2\times10^{-6}
\times5000\times0.01 = \qty{1.0e-4}{m/yr}$, \qty{0.1}{mm/yr}. Como $E$ está
em \unit{m/yr}, $A^m$ em $\unit{m}^{2m}$ e $S$ é adimensional, $K$ tem
unidade $\unit{m}^{1-2m}\,\unit{yr^{-1}}$: \unit{yr^{-1}} para $m = 0.5$.
\end{solution}

\begin{solution}{exo:b3:quantitative-geomorphology:4}
(a) Fixo: um leito duro torna o rio mais íngreme onde ele aflora, e o degrau
permanece com o leito. (b) Ancorado na falha: enquanto a falha continuar deslizando,
o rio fica mais íngreme onde a cruza. (c) Móvel: o soerguimento da
costa rebaixou o \omterm{def:b3:quantitative-geomorphology:knickpoint}{nível de base} de todos os rios de uma vez; os \omterm{def:b3:quantitative-geomorphology:knickpoint}{knickpoints}
estão migrando para montante e, como prevê o teorema da onda cinemática para $n
= 1$, situam-se sobre uma mesma curva de nível.
\end{solution}

\begin{solution}{exo:b3:quantitative-geomorphology:5}
$\theta = m/n = 0.5$. $S = (U/K)^{1/2}A^{-1/2} = (10^{-4}/10^{-9})^{1/2}
\times(10^7)^{-1/2} = 316\times3.16\times10^{-4} = 0.10$. Dobrar $U$
multiplica $k_s = (U/K)^{1/n}$ por $\sqrt2 = 1.41$; para $n = 1$, por 2. Quanto
maior $n$, menos a declividade responde ao soerguimento.
\end{solution}

\begin{solution}{exo:b3:quantitative-geomorphology:6}
Com $x$ contado a montante da foz, $\mathrm{d}x = -\mathrm{d}d$, logo $\chi(d)
= \int_d^L (A_0/c\,d'^{h})^{\theta}\,\mathrm{d}d' = A_0^{\theta} c^{-\theta}
\int_d^L d'^{-h\theta}\,\mathrm{d}d' = A_0^{\theta}c^{-\theta}(L^p -
d^p)/p$. Aqui $p = 0.19$, $A_0^{\theta}c^{-\theta} = 10^{2.7}/2^{0.45}$,
e $\chi = \qty{2935}{m}$.
\end{solution}

\begin{solution}{exo:b3:quantitative-geomorphology:7}
$c = 2\times10^{-6}\times(10^8)^{0.5} = \qty{0.02}{m/yr}$, \qty{20}{m/kyr};
\qty{5}{km} nessa velocidade levam \qty{250}{kyr}. A área, e com ela a
celeridade, diminui a montante, de modo que o \omterm{def:b3:quantitative-geomorphology:knickpoint}{knickpoint} desacelera à medida que sobe.
\end{solution}

\begin{solution}{exo:b3:quantitative-geomorphology:8}
$\bar e = 0.40\times600 = \qty{240}{m}$; os cumes se elevam em $0.80\times240 =
\qty{192}{m}$; a elevação média cai em $240 - 192 = \qty{48}{m}$.
\end{solution}

\begin{solution}{exo:b3:quantitative-geomorphology:9}
$h = h_0\ln(P_0/E)$: $0.5\ln 6 = \qty{0.90}{m}$ e $0.5\ln1.5 =
\qty{0.20}{m}$. A \qty{0.4}{mm/yr}, a erosão supera a produção mais rápida,
$P_0$: o solo é removido e a encosta passa a ser rocha nua, cujo rebaixamento
limitado pelo intemperismo.
\end{solution}

\begin{solution}{exo:b3:quantitative-geomorphology:10}
No \omterm{def:b3:quantitative-geomorphology:steady-profile}{estado estacionário}, $\mathrm{d}q/\mathrm{d}x = E$ e $q = 0$ no divisor, logo
$q = E x$. Então $\kappa S = E x\,(1 - S^2/S_c^2)$, uma quadrática cuja
raiz positiva é $S = [-\kappa + \sqrt{\kappa^2 + 4E^2x^2/S_c^2}\,]/(2Ex/S_c^2)$. Para
20, 50, 100 e \qty{200}{m}: 0.36, 0.68, 0.89 e 1.03, contra 0.4, 1.0,
2.0 e 4.0 para a difusão linear ($S = Ex/\kappa$). A colina tem uma crista arredondada
e, depois, flancos cujo gradiente tende a $S_c$: encostas quase planas, com
um relevo muito menor que o da parábola.
\end{solution}

\begin{solution}{exo:b3:quantitative-geomorphology:11}
Pelo teorema da onda cinemática, a mudança atinge um ponto de coordenada
$\chi$ após $\chi/(KA_0^m)$; a rede responde por completo quando atinge
a última cabeceira de canal, de maior $\chi$. $K A_0^m = 10^{-6}\times
1000 = \qty{e-3}{m/yr}$: o tronco responde em \qty{5.46}{Myr}, a rede
em \qty{7.24}{Myr}, da ordem de um milhão de anos, como se esperava.
\end{solution}

\begin{solution}{exo:b3:quantitative-geomorphology:12}
A montante da confluência, o gradiente estacionário do afluente é $S = (U/2K)A^{-\theta}$,
logo $\mathrm{d}z/\mathrm{d}\chi = (U/2K)A_0^{-\theta}$, a metade da do tronco.
O \omterm{def:b3:quantitative-geomorphology:knickpoint}{knickpoint} atinge a confluência quando o do tronco está em $\chi_j$; após um
novo intervalo de tempo $\Delta t$, o do tronco passou para $\chi_k$, com $\chi_k -
\chi_j = K A_0^m\,\Delta t$, enquanto o do afluente, com celeridade $2KA^m$,
passou para $\chi$, com $\chi - \chi_j = 2KA_0^m\,\Delta t = 2(\chi_k -
\chi_j)$.
\end{solution}

\begin{solution}{pb:b3:quantitative-geomorphology:1}
\textbf{1.} $4.5\times12000^{1.7} = \qty{3.87e7}{m^2}$, \qty{38.7}{km^2};
$4.5\times300^{1.7} = \qty{7.3e4}{m^2}$, \qty{0.073}{km^2}.

\textbf{2.} $\theta \approx -\log(0.007855/0.06641)/\log(37.08/0.388) =
0.928/1.980 = 0.47$.

\textbf{3.} Sim: com uma dispersão de cerca de \qty{10}{\percent} em doze pontos,
$0.466$ não difere significativamente de 0.5.

\textbf{4.} $0.06641\times\sqrt{388\,000} = \qty{41.4}{m}$;
$0.007855\times\sqrt{3.708\times10^7} = \qty{47.8}{m}$. A média geométrica,
\qty{45.1}{m}, estima $k_s$ para $\theta = 0.5$.

\textbf{5.} Um gradiente é a diferença de duas elevações, cada uma incerta
em decímetros a metros, dividida por uma distância curta; em trechos suaves, o
erro é uma parcela grande da diferença. $\chi$ integra, em vez de
derivar.

\textbf{6.} $p = 1 - 1.7\times0.5 = 0.15$; $\chi = (1000/\sqrt{4.5})\,
(12000^{0.15} - 9000^{0.15})/0.15 = \qty{543.1}{m}$.

\textbf{7.} $(193.4 - 170)/543.1 = 0.0431$; $k_{sn} = 0.0431\times A_0^{0.5}
= \qty{43.1}{m}$.

\textbf{8.} A estimativa por chi: ela se apoia em elevações ao longo de \qty{3}{km} de
rio, não em doze gradientes ruidosos.

\textbf{9.} $K = U/k_{sn} = 4.3\times10^{-5}/43.1 = \qty{1.00e-6}{yr^{-1}}$.

\textbf{10.} $10^{-6}\times\sqrt{3.87\times10^7}\times0.00691 =
\qty{4.30e-5}{m/yr}$, \qty{43}{m/Myr}: a foz acompanha o
levantamento.

\textbf{11.} Na foz, $10^{-6}\times6222 = \qty{6.2}{m/kyr}$; na
cachoeira, $d = \qty{7.40}{km}$, $A = \qty{1.70e7}{m^2}$ e
\qty{4.1}{m/kyr}.

\textbf{12.} $170 + 43\times(4.5)^{-0.5}(12000^{0.15} - 300^{0.15})/0.15 =
\qty{405}{m}$: um relevo de \qty{235}{m}.

\textbf{13.} Para $n = 1$, a equação é uma onda cinemática: pontos de ambos os
lados do degrau sobem $U$ por unidade de tempo, de modo que o degrau mantém sua altura,
os \qty{25}{m} da queda.

\textbf{14.} $\chi = (1000/\sqrt{4.5})(12000^{0.15} - 7400^{0.15})/0.15 =
\qty{899.4}{m}$.

\textbf{15.} $t = \chi/(KA_0^{0.5}) = 899.4/10^{-3} = \qty{899}{kyr}$.

\textbf{16.} Em $\chi = 899.4 + 300 = \qty{1199}{m}$:
$d^{0.15} = 12000^{0.15} - 1199\times0.15\times\sqrt{4.5}/1000$, $d =
\qty{6.25}{km}$, \qty{5.75}{km} a montante da foz.

\textbf{17.} Na confluência ($d = \qty{9.5}{km}$) $\chi = \qty{442.8}{m}$;
a cachoeira está onde o $\chi$ do afluente atinge \qty{899.4}{m}:
$s^{0.15} = 5000^{0.15} - (899.4 - 442.8)\times0.15\times\sqrt{1.5}/1000$,
$s = \qty{4.27}{km}$, \qty{0.73}{km} acima da confluência.

\textbf{18.} Suas confluências estão em $\chi = 1954$ e \qty{999}{m}, além de
\qty{899}{m}: o \omterm{def:b3:quantitative-geomorphology:knickpoint}{knickpoint} ainda não as ultrapassou, de modo que o tronco ali
ainda está em seu nível antigo.

\textbf{19.} $E = 5\times10^{-3}\times0.0086 = \qty{4.3e-5}{m/yr}$,
\qty{43}{m/Myr}: as colinas acima da cachoeira erodem à taxa de levantamento,
em \omterm{def:b3:quantitative-geomorphology:steady-profile}{estado estacionário}.

\textbf{20.} $150^2/(5\times10^{-3}) = \qty{4.5}{Myr}$, cinco vezes a idade da
cachoeira.

\textbf{21.} A base delas foi rebaixada \qty{25}{m} junto com o canal: suas
partes inferiores estão excessivamente íngremes, talvez deslizando, enquanto suas cristas ainda
erodem à taxa antiga. Uma onda de acentuação sobe por cada encosta; depois de alguns
milhões de anos elas retornam à mesma parábola, \qty{25}{m} mais baixa.

\textbf{22.} $t = \chi k_{sn}/(U A_0^{0.5})$ varia como $1/U$: de
$899/1.2 = 749$ a $899/0.8 = \qty{1124}{kyr}$.

\textbf{23.} Com $n = 2$, a celeridade cresce com o gradiente: a face íngreme
se deslocaria mais depressa do que os trechos ao redor e permaneceria nítida, mas
sua velocidade e a idade lida a partir dela seriam diferentes.

\textbf{24.} O \omterm{def:b3:quantitative-geomorphology:knickpoint}{nível de base} do rio Varden caiu
há \textbf{\qty[text-series-to-math=true]{900}{kyr}} (\qty{899}{kyr} pelo
cálculo).
\end{solution}
